% !TeX root = ../组合数学笔记.tex
% 由主文档通过 \input 载入；请编译主文档。
\setcounter{chapter}{9} % 保留原书章号。
\chapter{图、路径与树的计数}
\label{ch:graphs-trees}

道路、通信线路和相互关系都能用图表达。图论关注顶点之间如何连接，组合计数则进一步询问满足指定结构的图有多少个。首先固定顶点与边的含义，再讨论连通、遍历和树的编码。

\section{图、度数与关联计数}

\begin{definition}
有限\emph{简单无向图（simple undirected graph）} $G=(V,E)$ 由有限顶点集 $V$ 及其二元素子集构成的边集 $E$ 组成。边 $\{u,v\}$ 也记为 $uv$；此时称 $u,v$ 邻接。顶点 $v$ 的邻点集记为 $N(v)$，度数记为 $d(v)=|N(v)|$。
\end{definition}

简单图没有环和重边。子图可删除顶点和边；由 $U\subseteq V$ 导出的诱导子图 $G[U]$ 则保留原图中两个端点都在 $U$ 的全部边。两个图之间若存在保持邻接与非邻接的顶点双射，称它们\emph{同构（isomorphic）}。固定顶点集 $[n]$ 时，每对顶点可独立决定是否连边，故有 $2^{\binom n2}$ 个标号简单图；按同构分类不能一般地直接除以 $n!$。

\begin{theorem}[握手定理]
$\sum_{v\in V}d(v)=2|E|$。因此奇度顶点的个数为偶数。
\end{theorem}
\begin{proof}
计数关联对 $(v,e)$，其中顶点 $v$ 是边 $e$ 的端点。按顶点计数得左侧，按边计数每边有两个端点，得右侧。度数总和为偶数，故其中奇数项的个数也为偶数。
\end{proof}

\section{连通性、路径与 Euler 道}

相邻顶点依次连接成\emph{游走（walk）}；边不重复的游走称为\emph{迹（trail）}，顶点不重复的游走称为\emph{路径（path）}。除首尾重合外顶点均不同的闭游走称为\emph{圈（cycle）}，在简单图中长度至少为 $3$。任意两点之间存在路径的非空图称为连通图，其极大连通部分称为连通分支。

经过每条边恰一次的迹称为\emph{Euler 道（Euler trail）}；首尾相同者为 Euler 闭迹。以下假定至少有一条边，且所有非孤立顶点处于同一连通分支。

\begin{theorem}
图存在 Euler 闭迹当且仅当所有顶点度数均为偶数；存在首尾不同的 Euler 道当且仅当恰有两个奇度顶点，且它们正是两个端点。
\end{theorem}
\begin{proof}
沿迹经过一个非端点时，每次进入都配对一次离开，所以其度数必须为偶数；首尾不同则两端各有一条未配对边，给出必要性。反之，若度数全偶，从有边顶点出发不断沿未用边前进，有限性保证最终无边可走。除起点外，进入顶点后已用边数在本次进入时为奇数，总度数为偶数，所以至少还有一条未用边，故只能在起点结束。若尚有未用边，由连通性，已有闭迹上必有顶点与未用边相接；删去已有闭迹后各点度数仍为偶数，从该点在剩余图中重复构造闭迹，再插入原闭迹。每次至少增加一条边，最终得到所需闭迹。若恰有两个奇度顶点，临时在它们之间增加一条允许重复的辅助边，构造闭迹后删除该边，即得开迹。
\end{proof}

\emph{Hamilton 圈（Hamilton cycle）}要求每个顶点恰经过一次，其任务与 Euler 道不同。两个三角形只共享一个顶点的图所有度数为偶数，因而有 Euler 闭迹，却没有 Hamilton 圈：若要访问两个三角形的其余顶点，必须重复经过公共顶点。度数奇偶条件不能作为 Hamilton 性的判据。

\section{树的等价刻画与生成树}

\begin{definition}
连通且无圈的图称为\emph{树（tree）}，无圈图称为\emph{森林（forest）}。度数为 $1$ 的顶点称为叶子。包含原图全部顶点的树子图称为\emph{生成树（spanning tree）}。
\end{definition}

\begin{theorem}
对含 $n\ge1$ 个顶点的简单图，下列条件等价：图是树；任意两点之间恰有一条路径；图连通且有 $n-1$ 条边。
\end{theorem}
\begin{proof}
连通保证路径存在。若两点间有两条不同路径，从首次分离到再次相遇的部分构成圈，因此树中路径唯一；反之，圈上两点有两条路径，故路径唯一蕴含无圈与连通。对 $n\ge2$ 的树，取最长路径，其端点不能有路径外邻点，否则可延长；也不能邻接路径上非相邻点，否则出现圈，故端点是叶子。删去叶子及其一条边后仍为树，归纳得边数为 $n-1$。最后，任意连通图可反复删去圈上的一条边而不破坏连通，终得生成树；若原图本来只有 $n-1$ 条边，则没有额外边可删，原图就是树。
\end{proof}

由每个分支为树，含 $n$ 个顶点、$c$ 个分支的森林有 $n-c$ 条边。一棵树选定一个顶点作根后，可按到根的距离定义父子关系；根树并不自动具有孩子之间的左右次序，只有再指定次序才是有序树。

\section{Pr\"ufer 编码与 Cayley 公式}

对顶点集 $[n]$、$n\ge2$ 的树，重复 $n-2$ 次以下操作：删除当前最小标签叶子，并记录它唯一邻点的标签。所得长度 $n-2$ 的串称为\emph{Pr\"ufer 序列（Pr\"ufer sequence）}。一个标签 $v$ 恰出现 $d(v)-1$ 次，因为删除其他叶子时，每记录一次 $v$ 就删掉一条与它相连的边，直到它自己成为叶子或留在最后两点中。

反向从任意串 $(a_1,\ldots,a_{n-2})\in[n]^{n-2}$ 出发。保留全部标签；第 $i$ 步取尚保留且未出现在后缀 $a_i,\ldots,a_{n-2}$ 中的最小标签 $b_i$，连边 $b_ia_i$ 并删除 $b_i$；最后连剩下两点。此时剩余标签数总比后缀长度多二，故 $b_i$ 必存在；$a_i$ 尚保留且不同于 $b_i$。逆序看这一构造，每步把一个新叶子接入已有树，所以最终为树。编码时，当前不再出现于后缀的标签恰是当前叶子，故两过程逐步互逆。

\begin{theorem}[Cayley 公式]
对 $n\ge2$，顶点集 $[n]$ 上的树有 $n^{n-2}$ 棵；$n=1$ 时有一棵。
\end{theorem}
\begin{proof}
编码是树与 $[n]^{n-2}$ 之间的双射，后者有 $n^{n-2}$ 个元素。单顶点图单独计数。
\end{proof}

例如序列 $(3,3)$ 对应 $[4]$ 上以 $3$ 为中心的星形树。进一步，若 $n\ge2$ 且各 $d_i\ge1$、$\sum_i d_i=2n-2$，指定度数序列的树数为 $(n-2)!/\prod_i(d_i-1)!$，因为标签 $i$ 须在序列中出现 $d_i-1$ 次。

\section{矩阵树定理与计数的线性代数表示}

设 $A$ 为邻接矩阵，$D=\operatorname{diag}(d(1),\ldots,d(n))$，$L=D-A$ 称为 Laplace 矩阵。记 $L^{(r)}$ 为删去第 $r$ 行及第 $r$ 列的矩阵。

\begin{theorem}[矩阵树定理（matrix-tree theorem）]
对 $n\ge2$ 的有限连通简单图，其生成树数为 $\det L^{(r)}$，与 $r$ 的选择无关。
\end{theorem}
\begin{proof}
任意给每条边一个方向，构造 $n\times |E|$ 的关联矩阵 $B$：每列在起点处为 $1$、终点处为 $-1$，其余为零。直接计算得 $BB^{\mathsf T}=L$。删去第 $r$ 行，得矩阵 $\widetilde B$，故 $L^{(r)}=\widetilde B\widetilde B^{\mathsf T}$。由附录\ref{app:prerequisites}中的 Cauchy--Binet 公式，
\[
 \det L^{(r)}=\sum_{\substack{S\subseteq E\\|S|=n-1}}
                 \det(\widetilde B_S)^2.
\]
若 $S$ 含圈，沿圈取带符号的列和为零，行列式为零。若无圈，由森林边数公式知 $(V,S)$ 连通，因而为树。此时取一个不是 $r$ 的叶子，沿其只有一个非零元的行展开，并递归删去叶子与对应边，最后得到 $|\det\widetilde B_S|=1$。所以每棵生成树贡献 $1$，其他边集贡献零。
\end{proof}

以上证明的行列式预备可与讲义\cite{matrix-tree-notes}对照。三角形的约化 Laplace 矩阵为 $\begin{pmatrix}2&-1\\-1&2\end{pmatrix}$，行列式为 $3$，与删去任意一边得到的三棵生成树一致。

\section{习题}
\begin{enumerate}
  \item\label{ex:graphs-trees-1} 证明树中至少有两个叶子（顶点数至少为 $2$），并说明最大度数不超过 $2$ 的树是什么图。
  \item\label{ex:graphs-trees-2} 解码顶点集 $[5]$ 上的 Pr\"ufer 序列 $(2,2,4)$，列出边集与各点度数。
  \item\label{ex:graphs-trees-3} 求 $[5]$ 上度数分别为 $(3,2,1,1,1)$ 的标号树数。
  \item\label{ex:graphs-trees-4} 用删除一条边的直接计数与矩阵树定理分别求四边形的生成树数。
\end{enumerate}

